\section{AMC 8 2005}
        \begin{enumerate}
        	\item (\textit{AMC 8 2005}) Connie multiplies a number by 2 and gets 60 as her answer. However, she should
have divided the number by 2 to get the correct answer. What is the correct
answer?


 $\textbf{(A)}\ 7.5\qquad\textbf{(B)}\ 15\qquad\textbf{(C)}\ 30\qquad\textbf{(D)}\ 120\qquad\textbf{(E)}\ 240$


 
 \textbf{Answer: }B



\textbf{Solution}

If $x$ is the number, then $2x=60$ and $x=30$. Dividing the number by $2$ yields $\dfrac{30}{2} = \boxed{\textbf{(B)}\ 15}$.

 A second way to do it is to divide the number by $4$, as you multiplied by $2$ when you were supposed to divide by $2$. So, $\dfrac{60}{4} = \boxed{\textbf{(B)}\ 15}$.

 


	\item (\textit{AMC 8 2005}) Karl bought five folders from Pay-A-Lot at a cost of $\textdollar 2.50$ each.
Pay-A-Lot had a 20\%-off sale the following day. How much could
Karl have saved on the purchase by waiting a day?


 $\textbf{(A)}\ \textdollar 1.00  \qquad\textbf{(B)}\ \textdollar 2.00 \qquad\textbf{(C)}\ \textdollar 2.50\qquad\textbf{(D)}\ \textdollar 2.75 \qquad\textbf{(E)}\ \textdollar 5.00$


 
 \textbf{Answer: }C



\textbf{Solution}

Karl paid $5 \cdot 2.50 = \textdollar 12.50$. $20 \%$ of this cost that he saved is $12.50 \cdot .2 = \boxed{\textbf{(C)}\ \textdollar 2.50}$.

 



\textbf{Solution 2}

Each folder can also be $5/2$ dollars, and $20\%$ can be shown as $(1/5)$. We can multiply $(5/2) \cdot (1/5) = (1/2)$. $(1/2)$  is also $50$ cents or the amount of money that is saved after the $20\%$ discount. So each folder is $2.50-0.5 = \textdollar2$.Since Karl bought 5 folders all of the folders after the discount is $(5)(2) = 10$, and the money bought before the discount is $(5)(2.50) = \textdollar12.50$. To find the money Karl saves all we have to do is subtract $12.50 - 10 = 2.50$. Thus the answer is  $\boxed{\textbf{(C)}\ \textdollar 2.50}$.

 



\textbf{Solution 3}

Having a $20\%$ off sale is equivalent to a "Buy 4 get 1 free" sale. Thus, he would have saved the price of 1 folder, which is $\boxed{\textbf{(C)}\ \textdollar 2.50}$.

 -JeffersonJ

 


	\item (\textit{AMC 8 2005}) What is the minimum number of small squares that must be colored black so that a line of symmetry lies on the diagonal $\overline{BD}$ of square $ABCD$?

                    \begin{figure}[H]
                        \centering
                        \includegraphics[width=0.5\textwidth]{Problem/AMC_8_AMC_2005_3_Problem_Figure_1.png}
                    \end{figure}
                    


 $\textbf{(A)}\ 1\qquad\textbf{(B)}\ 2\qquad\textbf{(C)}\ 3\qquad\textbf{(D)}\ 4\qquad\textbf{(E)}\ 5$


 
 \textbf{Answer: }D



\textbf{Solution}

Rotating square $ABCD$ counterclockwise $45^\circ$ so that the line of symmetry $BD$ is a vertical line makes it easier to see that $\boxed{\textbf{(D)}\ 4}$ squares need to be colored to match its corresponding square.

 


	\item (\textit{AMC 8 2005}) A square and a triangle have equal perimeters. The lengths of the three sides of the triangle are 6.1 cm, 8.2 cm and 9.7 cm. What is the area of the square in square centimeters?


 $\textbf{(A)}\ 24\qquad\textbf{(B)}\ 25\qquad\textbf{(C)}\ 36\qquad\textbf{(D)}\ 48\qquad\textbf{(E)}\ 64$


 
 \textbf{Answer: }C



\textbf{Solution}

The perimeter of the triangle is $6.1+8.2+9.7=24$ cm. A square's perimeter is four times its sidelength, since all its sidelengths are equal. If the square's perimeter is $24$, the sidelength is $24/4=6$, and the area is $6^2=\boxed{\textbf{(C)}\ 36}$.

 


	\item (\textit{AMC 8 2005}) Soda is sold in packs of 6, 12 and 24 cans. What is the minimum number of packs needed to buy exactly 90 cans of soda?


 $\textbf{(A)}\ 4\qquad\textbf{(B)}\ 5\qquad\textbf{(C)}\ 6\qquad\textbf{(D)}\ 8\qquad\textbf{(E)}\ 15$


 
 \textbf{Answer: }B



\textbf{Solution}

Start by buying the largest packs first. After three $24$-packs, $90-3(24)=18$ cans are left. After one $12$-pack, $18-12=6$ cans are left. Then buy one more $6$-pack. The total number of packs is $3+1+1=\boxed{\textbf{(B)}\ 5}$.

 


	\item (\textit{AMC 8 2005}) Suppose $d$ is a digit. For how many values of $d$ is $2.00d5 > 2.005$?


 $\textbf{(A)}\ 0\qquad\textbf{(B)}\ 4\qquad\textbf{(C)}\ 5\qquad\textbf{(D)}\ 6\qquad\textbf{(E)}\ 10$


 
 \textbf{Answer: }C



\textbf{Solution}

We see that $2.0055$ works but $2.0045$ does not. The digit $d$ can be numbers $5$ through $9$, which is $\boxed{\textbf{(C)}\ 5}$ values.

 


	\item (\textit{AMC 8 2005}) Bill walks $\frac12$ mile south, then $\frac34$ mile east, and finally $\frac12$ mile south. How many miles is he, in a direct line, from his starting point?


 $\textbf{(A)}\ 1\qquad\textbf{(B)}\ 1\frac14\qquad\textbf{(C)}\ 1\frac12\qquad\textbf{(D)}\ 1\frac34\qquad\textbf{(E)}\ 2$


 
 \textbf{Answer: }B



\textbf{Solution}

We have a right-angle triangle with sides 3/4 and 1. We can divide the simple pythagorean theorem identity 3,4,5 by 4 to get 3/4, 1, 5/4. This shows us the answer is $\boxed{\textbf{B}\ 1\frac{1}{4}}$.

 


	\item (\textit{AMC 8 2005}) Suppose m and n are positive odd integers. Which of the following must also be an odd integer?


 $\textbf{(A)}\ m+3n\qquad\textbf{(B)}\ 3m-n\qquad\textbf{(C)}\ 3m^2 + 3n^2\qquad\textbf{(D)}\ (nm + 3)^2\qquad\textbf{(E)}\ 3mn$


 
 \textbf{Answer: }E



\textbf{Solution}

Assume \href{https://artofproblemsolving.com/wiki/index.php?title=WLOG}{WLOG} that $m$ and $n$ are both $1$. Plugging into each of the choices, we get $4, 2, 6, 16,$ and $3$. The only odd integer is $\boxed{\textbf{(E)}\ 3mn}$.

 


	\item (\textit{AMC 8 2005}) In quadrilateral $ABCD$, sides $\overline{AB}$ and $\overline{BC}$ both have length 10, sides $\overline{CD}$ and $\overline{DA}$ both have length 17, and the measure of angle $ADC$ is $60^\circ$. What is the length of diagonal $\overline{AC}$?

                    \begin{figure}[H]
                        \centering
                        \includegraphics[width=0.5\textwidth]{Problem/AMC_8_AMC_2005_9_Problem_Figure_1.png}
                    \end{figure}
                    


 $\textbf{(A)}\ 13.5\qquad\textbf{(B)}\ 14\qquad\textbf{(C)}\ 15.5\qquad\textbf{(D)}\ 17\qquad\textbf{(E)}\ 18.5$


 
 \textbf{Answer: }D



\textbf{Solution 1}

Because $\overline{AD} = \overline{CD}$, $\triangle ADC$ is an isosceles triangle with $\angle DAC = \angle DCA$. Angles in a triangle add up to $180^\circ$, and since $\angle ADC=60^\circ$, the other two angles are also $60^\circ$, and $\triangle ADC$ is an equilateral triangle. Therefore $\overline{AC}=\overline{DA}=\boxed{\textbf{(D)}\ 17}$.

 



\textbf{Solution 2}

We can divide $\overline{CD}$ in half and connect this point to A, dividing $\triangle ADC$ in half. This means the base will be $\frac{17}{2}$ and the hypotenuse will be 17. By using the Pythagorean's Theorem, we see that if the base and height are shared, the hypotenuse should be the same. This tells us that the length of $\overline{AC} = \boxed{(D) 17}$.

 



\textbf{Solution 3 (Easiest)}

Take an equilateral triangle with side length $17$. $\triangle ADC$ is congruent to this by $SAS$, hence it is equilateral. The answer is $\boxed{\textbf{(D)}17}$.

 


	\item (\textit{AMC 8 2005}) Joe had walked half way from home to school when he realized he was late. He ran the rest of the way to school. He ran 3 times as fast as he walked. Joe took 6 minutes to walk half way to school. How many minutes did it take Joe to get from home to school?


 $\textbf{(A)}\ 7\qquad\textbf{(B)}\ 7.3\qquad\textbf{(C)}\ 7.7\qquad\textbf{(D)}\ 8\qquad\textbf{(E)}\ 8.3$


 
 \textbf{Answer: }D



\textbf{Solution 1}

We can use the equation $d=rt$ where $d$ is the distance, $r$ is the rate, and $t$ is the time. The distances he ran and walked are equal, so $r_rt_r=r_wt_w$, where $r_r$ is the rate at which he ran, $t_r$ is the time for which he ran, $r_w$ is the rate at which he walked, and $t_w$ is the time for which he walked. Because he runs three times faster than he walks, $r_r=3r_w$. We want to find the time he ran, $t_r = \frac{r_wt_w}{t_r} = \frac{(r_w)(6)}{3r_w} = 2$ minutes. He traveled for a total of $6+2=\boxed{\textbf{(D)}\ 8}$ minutes.

 



\textbf{Solution 2}

We know that it took Joe $6$ minutes to walk halfway. If he ran $3$ times as fast as that, he went $\frac63 = 2$ minutes for the other half (as $t=\frac{d}{s}$ and if we multiply $s$ by $3$ we are essentially dividing $t$ by $3$). Therefore, we have $6 + 2 = 8$ minutes.

 



\textbf{Solution 3}

According to $d = rt$, and if we assume the distance from home to school is $1$ mile, then $\frac{1}{2} = 6r$, meaning $r = \frac{1}{12}\text{ mph}$. Then, given that he runs $3$ times the speed he walks, he runs at $\frac{3}{12}\text{ mph}$ or $\frac{1}{4}\text{ mph}$. Also, note that he only has $\frac{1}{2}$ mile left to reach school. Therefore, once again using $d = rt$, $\frac{1}{2} = \frac{1}{4}t$, and $t = 2$. Combining the given $6$ minutes and the deduced $2$ minutes, we reach $\boxed{8\text{ minutes}}$.

 -N.

 


	\item (\textit{AMC 8 2005}) The sales tax rate in Bergville is 6\%. During a sale at the Bergville Coat Closet, the price of a coat is discounted 20\% from its $90.00 price. Two clerks, Jack and Jill, calculate the bill independently. Jack brings up $90.00 and adds 6\% sales tax, then subtracts 20\% from this total. Jill brings up $90.00, subtracts 20% of the price, then adds 6% of the discounted price for sales tax. What is Jack's total minus Jill's total?


 $\textbf{(A)}\ - \textdollar1.06\qquad\textbf{(B)}\  - \textdollar 0.53\qquad\textbf{(C)}\ \textdollar 0\qquad\textbf{(D)}\ \textdollar 0.53\qquad\textbf{(E)}\ \textdollar 1.06$


 
 \textbf{Answer: }C



\textbf{Solution}

The price Jack rings up is $\textdollar{(90.00)(1.06)(0.80)}$. The price Jill rings up is $\textdollar{(90.00)(0.80)(1.06)}$. By the commutative property of multiplication, these quantities are the same, and the difference is $\boxed{\textbf{(C)}\ \textdollar 0}$.

 
Reminder: Commutative states abc = bac= acb =..... and is the same in any order. This only works for addition and multiplication.

 


	\item (\textit{AMC 8 2005}) Big Al, the ape, ate 100 bananas from May 1 through May 5. Each day he ate six more bananas than on the previous day. How many delicious bananas did Big Al eat on May 5?


 $\textbf{(A)}\ 20\qquad\textbf{(B)}\ 22\qquad\textbf{(C)}\ 30\qquad\textbf{(D)}\ 32\qquad\textbf{(E)}\ 34$


 
 \textbf{Answer: }D



\textbf{Solution}

There are $5$ days from May 1 to May 5. The number of bananas he eats each day is an arithmetic sequence. He eats $n$ bananas on May 5, and $n-4(6)=n-24$ bananas on May 1. The sum of this arithmetic sequence is equal to $100$.

 
 \begin{align*} \frac{n+n-24}{2} \cdot 5 &= 100\\ n-12&=20\\ n&=\boxed{\textbf{(D)}\ 32} \end{align*}


 



\textbf{Solution 2}

There are $5$ days from May 1 to May 5. If we set the first day as $n$, the second day can be expressed as $n+6$, the third as $n+12$, and so on, for five days. The sum $n+(n+6)+(n+12)+(n+18)+(n+24)$ is equal to $100$, as stated in the problem. We can write a very simple equation, that is: $5n+60=100$.Now all we do is just solve. $5n=40$, so Big Al eats $8$ bananas on the first day. The fifth day, $n+24$, is then $32$, which is your answer.

 



\textbf{Solution 3 (Lightning Quick)}

Simply realize that the middle term of the arithmetic sequence is the arithmetic mean of all terms, which is simply $\frac{100}{5}=20$. This means that the number of bananas the ape ate on May 5th is just $20+6*2=32$. Select $\boxed{D}$.

 You can do this quicker than you read it if you truly master arithmetic sequences.

 


	\item (\textit{AMC 8 2005}) The area of polygon $ABCDEF$ is 52 with $AB=8$, $BC=9$ and $FA=5$. What is $DE+EF$?

                    \begin{figure}[H]
                        \centering
                        \includegraphics[width=0.5\textwidth]{Problem/AMC_8_AMC_2005_13_Problem_Figure_1.png}
                    \end{figure}
                    


 $\textbf{(A)}\ 7\qquad\textbf{(B)}\ 8\qquad\textbf{(C)}\ 9\qquad\textbf{(D)}\ 10\qquad\textbf{(E)}\ 11$


 
 \textbf{Answer: }C



\textbf{Solution}

Notice that $AF + DE = BC$, so $DE=4$. Let $O$ be the intersection of the extensions of $AF$ and $DC$, which makes rectangle $ABCO$. The area of the polygon is the area of $FEDO$ subtracted from the area of $ABCO$.

 
 \[\text{Area} = 52 = 8 \cdot 9- EF \cdot 4\]

 Solving for the unknown, $EF=5$, therefore $DE+EF=4+5=\boxed{\textbf{(C)}\ 9}$.

 



\textbf{Solution 2}

Similar to solution 1, $AF + DE = BC$, so $DE=4$. Split the polygon into two rectangles by extending the $DE$ so it intersects $AB$. Let's say the length of $FE$ is equal to $x$. We can form the equation: 
 \[5x + 9(8-x) = 52\]. We see that $x = 5$, so we can add $5 + 4 = \boxed{\textbf{(C)}\ 9}$.

 \begin{verbatim}
- ArjunBhumula
\end{verbatim}




	\item (\textit{AMC 8 2005}) The Little Twelve Basketball League has two divisions, with six teams in each division. Each team plays each of the other teams in its own division twice and every team in the other division once. How many games are scheduled?


 $\textbf{(A)}\ 80\qquad\textbf{(B)}\ 96\qquad\textbf{(C)}\ 100\qquad\textbf{(D)}\ 108\qquad\textbf{(E)}\ 192$


 
 \textbf{Answer: }B



\textbf{Solution}

Within each division, there are $\binom {6}{2} = 15$ pairings, and each of these games happens twice. The same goes for the other division so that there are $4(15)=60$ games within their own divisions. The number of games between the two divisions is $(6)(6)=36$. Together there are $60+36=\boxed{\textbf{(B)}\ 96}$ conference games.

 



\textbf{Solution 2}

Each team plays 10 games in its own division and 6 games against teams in the other division. So each of the 12 teams plays 16 conference games. Because each game involves two teams, there are $\frac{12\times 16}{2}=\boxed{96}$ games scheduled. ~aopsav (Credit to AoPS Alcumus)

 



\textbf{Solution 3}

Suppose division 1 has teams $A, B, C, D, E, F$. Suppose division 2 has teams $G, H, I, J, K, L$. Then, for interdivisional games in Division 1, $A$ has 10 games scheduled. Because $A$ has already played $B$ twice in those 10 games, $B$ will not have to account for $A$, making its scheduled games 8. This continues until $F$ has 0 games scheduled. The same will happen for division 2 in interdivisional games. So, $2(10 + 8 + 6 + 4 + 2) = 60$. Now, for cross-divisional games, team $A$ will play $G, H, I, J, K,$ and $L$ once. Each team from division 1 plays 6 cross-divisional games, for a total of $6 \cdot 6 = 36$ games. Division 2 will NOT be counted like done for division 1 since we have already considered the games. This is the same concept applied for interdivisional games. (This is also the concept applied in the Handshake Formula.) Now, $30 + 30 + 36 = \boxed{96}$ games scheduled.

 


	\item (\textit{AMC 8 2005}) How many different isosceles triangles have integer side lengths and perimeter 23?


 $\textbf{(A)}\ 2\qquad\textbf{(B)}\ 4\qquad\textbf{(C)}\ 6\qquad\textbf{(D)}\ 9\qquad\textbf{(E)}\ 11$


 
 \textbf{Answer: }C



\textbf{Solution}

Let $b$ be the base of the isosceles triangles, and let $a$ be the lengths of the other legs. From this, $2a+b=23$ and $b=23-2a$. From triangle inequality, $2a>b$, then plug in the value from the previous equation to get $2a>23-2a$ or $a>5.75$. The maximum value of $a$ occurs when $b=1$, in which from the first equation $a=11$. Thus, $a$ can have integer side lengths from $6$ to $11$, and there are $\boxed{\textbf{(C)}\ 6}$ triangles.

 


	\item (\textit{AMC 8 2005}) A five-legged Martian has a drawer full of socks, each of which is red, white or blue, and there are at least five socks of each color. The Martian pulls out one sock at a time without looking. How many socks must the Martian remove from the drawer to be certain there will be 5 socks of the same color?


 $\textbf{(A)}\ 6\qquad\textbf{(B)}\ 9\qquad\textbf{(C)}\ 12\qquad\textbf{(D)}\ 13\qquad\textbf{(E)}\ 15$


 
 \textbf{Answer: }D



\textbf{Solution}

The Martian can pull out $12$ socks, $4$ of each color, without having $5$ of the same kind yet. However, the next one he pulls out must be the fifth of one of the colors so he must remove $\boxed{\textbf{(D)}\ 13}$ socks.

 


	\item (\textit{AMC 8 2005}) The results of a cross-country team's training run are graphed below. Which student has the greatest average speed?

                    \begin{figure}[H]
                        \centering
                        \includegraphics[width=0.5\textwidth]{Problem/AMC_8_AMC_2005_17_Problem_Figure_1.png}
                    \end{figure}
                    


 $\textbf{(A)}\ \text{Angela}\qquad\textbf{(B)}\ \text{Briana}\qquad\textbf{(C)}\ \text{Carla}\qquad\textbf{(D)}\ \text{Debra}\qquad\textbf{(E)}\ \text{Evelyn}$


 
 \textbf{Answer: }E



\textbf{Solution}

Average speed is distance over time, or the slope of the line through the point and the origin. $\boxed{\textbf{(E)}\ \text{Evelyn}}$ has the steepest line, and runs the greatest distance for the shortest amount of time.

 


	\item (\textit{AMC 8 2005}) How many three-digit numbers are divisible by 13?


 $\textbf{(A)}\ 7\qquad\textbf{(B)}\ 67\qquad\textbf{(C)}\ 69\qquad\textbf{(D)}\ 76\qquad\textbf{(E)}\ 77$


 
 \textbf{Answer: }C



\textbf{Solution 1}

Let $k$ be any positive integer so that $13k$ is a multiple of $13$. For the smallest three-digit number, $13k>100$ and $k>\frac{100}{13} \approx 7.7$. For the greatest three-digit number, $13k<999$ and $k<\frac{999}{13} \approx 76.8$. The number $k$ can range from $8$ to $76$ so there are $\boxed{\textbf{(C)}\ 69}$ three-digit numbers.

 


 



\textbf{Solution 2}

3 digits number start from 100 and ends at 999.1 to 100 total multiple of $13 =\frac{100}{13}$ 

 1 to 999 total multple of $13 =\frac{999}{13}$ 

 Hence $\frac{999}{13}-\frac{100}{13}$=76-7=69

 


	\item (\textit{AMC 8 2005}) What is the perimeter of trapezoid $ABCD$?


 
                    \begin{figure}[H]
                        \centering
                        \includegraphics[width=0.5\textwidth]{Problem/AMC_8_AMC_2005_19_Problem_Figure_1.png}
                    \end{figure}
                    


 $\textbf{(A)}\ 180\qquad\textbf{(B)}\ 188\qquad\textbf{(C)}\ 196\qquad\textbf{(D)}\ 200\qquad\textbf{(E)}\ 204$


 
 \textbf{Answer: }A



\textbf{Solution}

Draw altitudes from $B$ and $C$ to base $AD$ to create a rectangle and two right triangles. The side opposite $BC$ is equal to $50$. The bases of the right triangles can be found using Pythagorean or special triangles to be $18$ and $7$. Add it together to get $AD=18+50+7=75$. The perimeter is $75+30+50+25=\boxed{\textbf{(A)}\ 180}$.

 


	\item (\textit{AMC 8 2005}) Alice and Bob play a game involving a circle whose circumference is divided by 12 equally-spaced points. The points are numbered clockwise, from 1 to 12. Both start on point 12. Alice moves clockwise and Bob, counterclockwise.
In a turn of the game, Alice moves 5 points clockwise and Bob moves 9 points counterclockwise. The game ends when they stop on the same point. How many turns will this take?


 $\textbf{(A)}\ 6\qquad\textbf{(B)}\ 8\qquad\textbf{(C)}\ 12\qquad\textbf{(D)}\ 14\qquad\textbf{(E)}\ 24$


 
 \textbf{Answer: }A



\textbf{Solution}

Alice moves $5k$ steps and Bob moves $9k$ steps, where $k$ is the turn they are on. Alice and Bob coincide when the number of steps they move collectively, $14k$, is a multiple of $12$. Since this number must be a multiple of $12$, as stated in the previous sentence, $14$ has a factor $2$, $k$ must have a factor of $6$. The smallest number of turns that is a multiple of $6$ is $\boxed{\textbf{(A)}\ 6}$.

 Another solution (Modular Arithmetic) Step 1: Write their positions after 𝑘k turns.

 Alice: moves 5 steps clockwise each turn.Position after 𝑘k turns:

 𝐴(𝑘)=12+5𝑘(mod12)A(k)=12+5k(mod12)

 Bob: moves 9 steps counterclockwise each turn.Moving 9 steps counterclockwise is the same as moving +3 clockwise (since 12-9=312-9=3).Position after 𝑘k turns:

 𝐵(𝑘)=12+3𝑘(mod12)B(k)=12+3k(mod12)Step 2: Find when they meet.

 They meet when:

 𝐴(𝑘)\equiv𝐵(𝑘)(mod12)A(k)\equivB(k)(mod12)

 So:

 12+5𝑘\equiv12+3𝑘(mod12)12+5k\equiv12+3k(mod12)

 Cancel the 12’s:

 5𝑘\equiv3𝑘(mod12)5k\equiv3k(mod12)2𝑘\equiv0(mod12)2k\equiv0(mod12)Step 3: Solve for 𝑘k.

 That means 2𝑘2k is divisible by 12:

 2𝑘\equiv0(mod12)\Rightarrow𝑘\equiv0(mod6)2k\equiv0(mod12)\Rightarrowk\equiv0(mod6)

 So the smallest positive 𝑘k is:

 𝑘=6(True)k=6

 (True) They meet after 6 turns → Answer: (A) 6  - timi821

 


	\item (\textit{AMC 8 2005}) How many distinct triangles can be drawn using three of the dots below as vertices?


 
                    \begin{figure}[H]
                        \centering
                        
\includegraphics[width=0.5\textwidth]{Problem/AMC_8_AMC_2005_21_Problem_Figure_1.png}
                    \end{figure}
                    


 $\textbf{(A)}\ 9\qquad\textbf{(B)}\ 12\qquad\textbf{(C)}\ 18\qquad\textbf{(D)}\ 20\qquad\textbf{(E)}\ 24$


 
 \textbf{Answer: }C



\textbf{Solution 1}

The number of ways to choose three points to make a triangle is $\binom 63 = 20$. However, two* of these are a straight line so we subtract $2$ to get $\boxed{\textbf{(C)}\ 18}$.

 
                \begin{itemize}
                    \item Note:  We are assuming that there are no degenerate triangles in this problem, and that is why we subtract two.
                \end{itemize}


                



\textbf{Solution 2}

\textbf{Case 1: One vertex is on the top 3 points}

 Here, there is $\binom 31 =3$ ways to choose the vertex on the top and $\binom 32 =3$ ways the choose the $2$ on the bottom, so there is $3 \cdot 3=9$ triangles.

 \textbf{Case 2: One vertex is on the bottom 3 points}

 By symmetry, there is $9$ triangles.

 Otherwise, the triangle is degenerate. 

 Our final answer is $9+9=\boxed{\text{(C) 18}}$ ~\href{https://artofproblemsolving.com/wiki/index.php?title=Bobbert&action=edit&redlink=1}{Bobbert}

 


	\item (\textit{AMC 8 2005}) A company sells detergent in three different sized boxes: small (S), medium (M) and large (L). The medium size costs 50\% more than the small size and contains 20\% less detergent than the large size. The large size contains twice as much detergent as the small size and costs 30\% more than the medium size. Rank the three sizes from best to worst buy.


 $\textbf{(A)}\ \text{SML}\qquad\textbf{(B)}\ \text{LMS}\qquad\textbf{(C)}\ \text{MSL}\qquad\textbf{(D)}\ \text{LSM}\qquad\textbf{(E)}\ \text{MLS}$


 
 \textbf{Answer: }E



\textbf{Solution}

WLOG (Without Loss of Generality), suppose the small size costs $\textdollar 1$ and has $5$ oz. The medium size then costs $\textdollar 1.50$ and has $8$ oz. The large size costs $\textdollar 1.95$ and has 10 oz. The small, medium, and large sizes cost, respectively, $0.200,$0.188, and S0.195 per oz. The sizes from best to worst buy are $\boxed{\textbf{(E)}\ \text{MLS}}$.

 


	\item (\textit{AMC 8 2005}) Isosceles right triangle $ABC$ encloses a semicircle of area $2\pi$. The circle has its center $O$ on hypotenuse $\overline{AB}$ and is tangent to sides $\overline{AC}$ and $\overline{BC}$. What is the area of triangle $ABC$?


 
                    \begin{figure}[H]
                        \centering
                        \includegraphics[width=0.5\textwidth]{Problem/AMC_8_AMC_2005_23_Problem_Figure_1.png}
                    \end{figure}
                    


 $\textbf{(A)}\ 6\qquad\textbf{(B)}\ 8\qquad\textbf{(C)}\ 3\pi\qquad\textbf{(D)}\ 10\qquad\textbf{(E)}\ 4\pi$


 
 \textbf{Answer: }B



\textbf{Solution}

First, we notice half a square so first let's create a square. Once we have a square, we will have a full circle. This circle has a diameter of 4 which will be the side of the square. The area would be $4\cdot 4 = 16.$ Divide 16 by 2 to get the original shape and you get $\boxed{8}$

 



\textbf{Solution 2 (Super Fast)}

Well, the question states that it is an Isosceles triangle, so 2 of the side lengths (not the hypotenuse) must be the same. Let us call them x, and the formula for finding the area of a triangle is 1/2(a)(b), if we implement this we get 1/2(x)(x) = x\textasciicircum{}2/2. The only answer which satisfies this condition is B therefore we have solved the question. We can double check as 8x2=16 and the root of 16 = 4.

 $\boxed{\textbf{(B)}\ 8 }$.

 


	\item (\textit{AMC 8 2005}) A certain calculator has only two keys [+1] and [x2]. When you press one of the keys, the calculator automatically  displays the result. For instance, if the calculator originally displayed "9" and you pressed [+1], it would display "10." If you then pressed [x2], it would display "20." Starting with the display "1," what is the fewest number of keystrokes you would need to reach "200"?


 $\textbf{(A)}\ 8\qquad\textbf{(B)}\ 9\qquad\textbf{(C)}\ 10\qquad\textbf{(D)}\ 11\qquad\textbf{(E)}\ 12$


 
 \textbf{Answer: }B



\textbf{Solution 1 (Now rigorous)}

We can start at $200$ and work our way down to $1$. We want to press the button that multiplies by $2$ the most, but since we are going down instead of up, we divide by $2$ instead. If we come across an odd number, then we will subtract that number by $1$. Notice

 \begin{verbatim}
$200 \div 2 = 100$,  $100 \div 2 = 50$,  $50 \div 2 = 25$,$25-1 = 24$,(Cannot divide by 2)  $24 \div 2 = 12$,  $12 \div 2 = 6$,  $6 \div 2 = 3$,  $3-1 = 2$,(Cannot divide by 2)   $2 \div 2 = 1$.   
\end{verbatim}

Since we've reached $1$, it's clear that the answer should be $\boxed{\textbf{(B)}\ 9}$- $\boxed{\textbf{Javapost}}$. Because we only subtracted $1$ when we had to, this is optimal. ~Roy2020 minor edits by IncredibleCheetah53

 



\textbf{Solution 2 (duodecimal)}

We make two key observations. First, pressing [x2] appends a $0$ to the end of a number's binary representation. Second, pressing [x2] and then pressing [+1] appends to the end of a number's binary representation.

 The base-ten number $200$ is represented as $11001000_2$ in binary. Therefore, the five $0$s contribute $5$ button presses. Similarly, each of the three $1$s each contribute $2$ button presses, although since we initially started with a $1$ we can ignore one of the $1$s. Thus, the answer is $5 + 2 \cdot 2 = \boxed{\textbf{(B)}\ 9}$.

 


	\item (\textit{AMC 8 2005}) A square with side length 2.0 and a circle share the same center. The total area of the regions that are inside the circle and outside the square is equal to the total area of the regions that are outside the circle and inside the square. What is the radius of the circle?


 
                    \begin{figure}[H]
                        \centering
                        \includegraphics[width=0.5\textwidth]{Problem/AMC_8_AMC_2005_25_Problem_Figure_1.png}
                    \end{figure}
                    
$\textbf{(A)}\ \frac{2}{\sqrt{\pi}} \qquad \textbf{(B)}\ \frac{1+\sqrt{2}}{2} \qquad \textbf{(C)}\ \frac{3}{2} \qquad \textbf{(D)}\ \sqrt{3} \qquad \textbf{(E)}\ \sqrt{\pi}$


 
 \textbf{Answer: }A



\textbf{Solution 1}

Let the region within the circle and square be $a$. In other words, it is the area inside the circle $\textbf{and}$ the square.  Let $r$ be the radius. We know that the area of the circle minus $a$ is equal to the area of the square, minus $a$ . 

 We get:

 $\pi r^2 -a=4-a$

 $r^2=\frac{4}{\pi}$

 $r=\frac{2}{\sqrt{\pi}}$

 So the answer is $\boxed{\textbf{(A)}\ \frac{2}{\sqrt{\pi}}}$.

 



\textbf{Solution 2}

We realize that since the areas of the regions outside of the circle and the square are equal to each other, the area of the circle must be equal to the area of the square.

 $\pi r^2=4$

 $r^2=\frac{4}{\pi}$

 $r=\frac{2}{\sqrt{\pi}}$

 So the answer is $\boxed{\textbf{(A)}\ \frac{2}{\sqrt{\pi}}}$.

 


\end{enumerate}